% 本文件由 LaTeX 排版 Agent 根据有效 translation.json 自动生成。
% author=Tertullian; work_id=0322; title=论祈祷

\chapter{论祈祷}

% structure: p0001 source=h1 target=omitted reason=page-title-already-rendered-by-parent

% structure: p0002 source=h2 target=section reason=relative-heading-rank; base=h2

\section{第一章 总论}

天主之神、天主之圣言与天主之理性——圣言之理性、理性与圣言之神——即我们的主耶稣基督，祂兼具二者——为我们新约的门徒规定了一种新的祈祷形式；因为在这事上也必须将新酒装在新皮囊里，将新布补在新衣服上。此外，往昔的一切，或已完全改变，如割礼；或已补充，如律法的其余部分；或已应验，如预言；或已成全，如信德本身。因为天主的新恩宠藉着引入福音——这福音是往昔整个旧制度的抹除者——将一切从属肉体的更新为属神的；在其中，我们的主耶稣基督被证明是天主之神、天主之圣言和天主之理性：祂藉以行大能的是神，祂藉以教导的是圣言，祂藉以降临的是理性。因此，基督所编排的祈祷包含三部分：在言词上，藉此\Emph{祈祷}被说出；在神上，单藉此祈祷才能生效。就连若翰也教导过他的门徒祈祷，但若翰的一切作为都是为基督预备根基，直到\Quote{祂增长了}——正如同一若翰曾预告\Quote{必须}使\Quote{祂增长，而他自己衰微}\BibleRef{若 3:30}——前驱的全部工作连同其神本身都转移到了主身上。因此，若翰教导祈祷的言辞形式并未流传，因为属地的已让位给属天的。若翰说：\Quote{那从地上来的，讲论属地的事；那从天上来的，讲论祂所看见的事。}\BibleRef{若 3:31}那么，主基督的——正如这种祈祷方式——难道不是属天的吗？因此，蒙福的\Emph{弟兄们}，让我们思量祂属天的智慧：首先，关于秘密祈祷的诫命，祂藉此要求人的信德，使他确信全能天主的视听闻达于屋顶之下，甚至深入内室；并要求信德上的谦逊，使其只向祂献上宗教的敬意，因为人相信祂无所不见、无所不闻。此外，既然智慧继之以如下的诫命，愿它同样归属于信德及信德的谦逊，使我们不认为必须以繁文缛词来接近主，因为我们确信祂不待祈求就已为祂自己的人预见。然而，正是这简短——且让这构成智慧的第三等级——是由伟大而蒙福的诠释之实质所支撑的，其意义之广博如同其言辞之凝练。因为它不仅囊括了祈祷的特殊职责，无论是尊敬天主还是为人祈求，而且几乎涵盖了主的一切言论、祂训诲的一切记录；事实上，在祈祷中包含了全部福音的纲要。\par

% structure: p0004 source=h2 target=section reason=relative-heading-rank; base=h2

\section{第二章 第一段}

祈祷以对天主的见证和信德的赏报开始，当我们说：\Quote{我们在天上的父}时；因为（这样说时）我们既向天主祈祷，又赞许信德，而这称呼正是信德的赏报。经上记载：\Quote{凡信祂的人，祂赐给他们权能，得以称为天主的儿女。}\BibleRef{若 1:12}然而，我们的主经常向我们宣示天主为父；甚至颁布诫命\Quote{不要称呼地上的人为父，惟独要称我们在天上的父}；\BibleRef{玛 23:9}因此，如此祈祷时，我们也在遵行这诫命。那些认识自己父亲的人是有福的！这正是对以色列的责备，圣神为之作证天地，说：\Quote{我养育了儿女，他们却不认识我。}\BibleRef{依 1:2}再者，当说\Quote{父}时，我们也称祂为\Quote{天主}。这称呼既关乎孝爱之责，也关乎权能。此外，在父中，子也被呼求；因为祂说：\Quote{我与父原是一体。}\BibleRef{若 10:30}甚至连我们的母亲教会也没有被忽略，也就是说，若在父与子中承认了母亲，父亲和儿子的名号正是从她而来。因此，用一个总称或一个词，我们既尊崇天主连同祂的所有，又谨记诫命，并给那些忘记了他们父亲的人打上记号。\par

% structure: p0006 source=h2 target=section reason=relative-heading-rank; base=h2

\section{第三章 第二段}

\Quote{天主父}这名从未向任何人启示过。甚至\Person{梅瑟}，他曾就这一点询问过天主，也只听到一个不同的名字。\BibleRef{出 3:13} 但为我们，这名字却藉圣子启示了，因为圣子如今就是圣父的新名。祂说：\Quote{我以圣父的名义而来}\BibleRef{若 5:43}；又说：\Quote{父啊，光荣你的名}\BibleRef{若 12:28}；更明确地说：\Quote{我已将你的名显示给人们}\BibleRef{若 17:6}。因此，我们祈求那\Emph{名}得以\Quote{被尊为圣}。并非因为人应当\Emph{祝愿}天主\Emph{安好}，仿佛另有人能祝愿祂，或除非我们如此祝愿，祂便会受损。显然，因着人人永远应纪念祂的恩惠，天主在每一处、每一时受\Emph{赞美}本是普遍合宜的。但这祈求也起到祝福的作用。否则，天主的圣名何时不\Quote{圣}，不藉着祂自己成为\Quote{圣}呢？祂亲自圣化万物——那围绕祂的天使不断呼喊：\Quote{圣！圣！圣！}因此，我们这些争取成为天使的人，若堪得此恩，从今世便开始预习那将来要向天主颂唱的诗歌和未来荣耀的职分。以上是为天主的荣耀。另一方面，当我们说\Quote{愿你的名受显扬}时，这也是为我们自己的祈求：愿这名在\Emph{我们}内——我们这些在祂内的人——受显扬，也在所有其他等待天主恩宠的人内受显扬。\BibleRef{依 30:18} 使我们也能遵守这条规诫：\Quote{为所有人祈祷}，\BibleRef{弟前 2:1} 甚至为我们个人的仇敌祈祷。\BibleRef{玛 5:44} 因此，我们欲言又止，不说\Quote{愿它\Emph{在我们内}受显扬}，而是说\Quote{愿它受显扬\Emph{在所有人内}}。\par

% structure: p0008 source=h2 target=section reason=relative-heading-rank; base=h2

\section{第四章 第三祈求}

照此模式，我们加上：\Quote{愿祢的旨意奉行在天上及地上}；并非有某种力量阻止天主旨意的实现，以致我们祈求祂的旨意成功实现；而是我们祈求祂的旨意奉行在\Emph{一切内}。因为藉着对\Emph{肉体}和\Emph{精神}的比喻解释，\Emph{我们}就是\Quote{天}和\Quote{地}；即使按字面理解，这祈求的意义也一样：愿\Emph{在我们内}，天主的旨意得以奉行在地上，从而使之也能奉行在天上。此外，天主\Emph{究竟}愿意什么，不就是我们按祂的训诲行事吗？因此，我们祈求祂赐予我们祂旨意的实质和履行它的能力，使我们能在天上和地上得救；因为祂旨意的总汇就是祂所收纳之人的救恩。还有那一个天主的旨意，即主在宣讲、行事、忍受中所成就的：因为如果祂亲自宣称祂所行的不是自己的，而是圣父的旨意，那么无疑祂所行的那些事\Emph{就是}圣父的旨意；\BibleRef{若 6:38} 这些事作为榜样，如今激励我们也去宣讲、行事，甚至忍受至死。我们\Emph{需要}天主的旨意，才能履行这些职责。再者，当我们说\Quote{愿祢的旨意奉行}时，我们甚至是祝愿我们自己安好，因为天主的旨意中毫无\Emph{邪恶}；即使按各人应得，可能另有别的加于我们。因此，藉此表达我们提醒自己忍耐。主也曾以自己肉身的真实受苦，向我们显示肉身的软弱，祂说：\Quote{父啊，请移开这杯}；又回想到自己，\Emph{加上}：\Quote{但不要照我的意愿，而要照祢的旨意成就。}\BibleRef{路 22:42} 祂\Emph{就是}圣父的旨意和能力；然而，为显示应得的忍耐，祂将自己交付\Emph{于}圣父的旨意。\par

% structure: p0010 source=h2 target=section reason=relative-heading-rank; base=h2

\section{第五章 第四祈求}

\Quote{愿祢的国来临}也指向\Quote{愿祢的旨意奉行}所指的——即\Emph{在我们内}。因为什么时候天主不统治呢？万王之心都在祂手中。\BibleRef{箴 21:1} 但我们为自己所愿的，我们为祂占卜，并且将我们预期\Emph{从}祂而来的归于\Emph{祂}。因此，如果主的国的\Emph{显现}属于天主的旨意和我们的热切期待，那么当所求的天国即将来临，趋向世代的终结时，何以有人祈求延续世代呢？我们的愿望是加速我们为王，而不是延长我们的奴役。即使祈祷中没有规定我们要祈求天国的来临，我们也会自发地发出那呼声，加速我们希望的实现。祭台下殉道者的灵魂嫉妒地向主呼喊说：\Quote{圣洁真实的主啊，你不审判住在地上的人，给我们伸流血的冤，要到几时呢？}\BibleRef{默 6:10} 当然，他们的伸冤是根据世代的终结。不，主，愿你的国速速来临——基督徒的祈祷，外邦人的混乱，天使的欢腾，我们为此受苦，不，我们为此祈祷！\par

% structure: p0012 source=h2 target=section reason=relative-heading-rank; base=h2

\section{第六章 第五祈求}

然而，上智安排祷词的次序何其巧妙，使它在属天之事——即天主的\Quote{名}、天主的\Quote{旨意}和天主的\Quote{国}——\Emph{之后}，也为地上的需要留出祈求的空间！因为主曾颁布谕令说：\BibleQuote{你们先寻求天国，这一切自会加给你们。}（\BibleRef{玛 6:33}）尽管我们也可以从灵性\Emph{角度}理解\Quote{求祢今天赐给我们日用的食粮}。因为\Emph{基督}是我们的食粮；因为基督是生命，而食粮就是生命。祂说：\BibleQuote{我是生命之粮}（\BibleRef{若 6:35}）；稍前又说：\BibleQuote{这食粮是降自永生活天主的圣言。}（\BibleRef{若 6:33}）再者，\Emph{我们发现}祂的身体也被视为食粮：\BibleQuote{这是我的身体。}（\BibleRef{玛 26:26}）因此，当我们祈求\Quote{日用食粮}时，我们是在祈求常在其内，并与祂的身体不可分离。然而，此词也可按肉性意义理解，但使用时不可不兼顾灵修操练的神圣纪念；因为主命令我们祈求食粮，这原是信徒唯一必需的饮食；\BibleQuote{其余一切，都是外邦人所求的}（\BibleRef{玛 6:32}）。祂同样藉比喻教导这教训，反复在比喻中论及，祂说：\Quote{父亲岂会从儿子手中夺去食粮，丢给狗吃呢？}又说：\Quote{父亲岂会在儿子求食粮时，给他一块石头呢？}祂\Emph{如此}表明儿子应期望从父亲得到什么。甚至那个半夜叩门者所求的也是\Emph{食粮}（\BibleRef{路 11:5}）。再者，祂恰当地加上\Quote{今天}，因为祂先前曾说：\Quote{不要为明天忧虑吃什么。}祂还为此配以比喻，讲论一个人思量扩建粮仓以储藏将来的收成，并期待长久的安逸；但那夜他便死了。（\BibleRef{路 12:16}）\par

% structure: p0014 source=h2 target=section reason=relative-heading-rank; base=h2

\section{第七章 第六项祈求}

在默想天主的慷慨之后，理当也向祂的仁慈祈求。因为食物若果真将我们如同待宰的公牛般\Emph{交付}于它们，又有何益？主自知是唯一无罪者，因而教导我们祈求\BibleQuote{赦免我们的罪债。}求赦免即是完全的告明；因为求赦免者全然承认自己的罪过。如此，悔改被表明为天主所悦纳，祂喜爱悔改胜于罪人的死亡。此外，在圣经中，\Quote{债务}是\Quote{罪过}的象徵；因为它同样应受审判的判决，并被判决所追讨；除非追讨被豁免，否则无法逃避正义的追讨，正如那比喻中主人豁免了那奴仆的债务（\BibleRef{玛 18:21}）；因为整个比喻的用意正在于此。事实上，那仆人被主人释放后，并未同样宽免自己的债务人；因此被主控诉，被交给刑役，直到还清最后一分钱——即每一罪过，无论多小——这正符合我们的宣认：\BibleQuote{我们也宽免我们的债务人；}的确，在其他地方，与这祷词形式一致，祂说：\BibleQuote{你们宽免，也必蒙宽免。}（\BibleRef{路 6:37}）当伯多禄问是否应七次宽免弟兄时，祂说：\Quote{不，七十个七次}（\BibleRef{玛 18:21}），为要改良法律；因为在《创世纪》中，加音被报应\Quote{七次}，而拉默客则\Quote{七十个七次}。\par

% structure: p0016 source=h2 target=section reason=relative-heading-rank; base=h2

\section{第八章 第七项即末项祈求}

为使这如此简短的祷词臻于圆满，祂加了——使我们不仅祈求赦免，更祈求完全避免罪过——\Quote{不要让我们陷入诱惑：}即，求祢不容许我们被那诱惑者所诱；但切莫认为\Emph{天主}会诱惑人，好像祂不知某人的信德，或急于推翻它。软弱和恶意是魔鬼的特性。因为天主甚至命令亚巴郎祭献自己的儿子，并非为诱惑，而是为考验他的信德；为藉他树立一个祂日后要颁布的诫命的榜样，即不可将任何爱慕的抵押看得比天主更珍贵。主自己受魔鬼诱惑时，已指明谁主管并创始诱惑。祂用后来的话证实这段话，说：\Quote{你们祈祷，免得陷于诱惑}；然而他们却因放弃祈祷而陷入睡眠，背弃了主，故\Emph{确实}受了诱惑。因此，末项祈求与\Quote{不要让我们陷入诱惑}的\Emph{含义}一致；因为其含义是：\Quote{但救我们脱离那恶者。}\par

% structure: p0018 source=h2 target=section reason=relative-heading-rank; base=h2

\section{第九章 总结}

在这寥寥数语的总结中，先知书、福音书及宗徒们的多少言论，主自己的多少讲论、榜样和比喻，被触及！多少义务同时履行！在\Quote{父}中，是天主的荣耀；在\Quote{名}中，是信德的见证；在\Quote{旨意}中，是顺从的奉献；在\Quote{国}中，是望德的纪念；在\Quote{粮}中，是生命的祈求；在祈求\Quote{宽恕}中，是罪债的完全承认；在祈求\Quote{护佑}中，是对诱惑的焦虑畏惧。有何希奇？唯天主能教导祂愿意如何被祈祷。因此，这由祂亲自设立、并由祂的圣神在从神圣口中发出时赋予生气的祈祷礼仪，凭其特权升入天堂，将圣子所教导的呈于圣父。\par

% structure: p0020 source=h2 target=section reason=relative-heading-rank; base=h2

\section{第十章 我们可在主祷文上加添自己的祈祷}

然而，主既是人类需要的预见者，在传授祂的祈祷规则之后，又分别说：\Quote{你们求，就必得着}；并且\Emph{因为}有按各人情况而作的祈祷，我们额外的需要就有权——在开始以合法和惯常的祈祷作为基础之后——建造一个外在的祈求上层建筑，但仍要记念\Emph{师傅的}训诫。\par

% structure: p0022 source=h2 target=section reason=relative-heading-rank; base=h2

\section{第十一章 祈祷圣父时，不可对弟兄动怒}

为使我们不致于远离天主的耳朵如同远离祂的诫命一样，记念祂的诫命为我们的祈祷铺平了一条通往天堂的道路；这些\Emph{诫命}中首要的是，我们不可上到天主的祭台前，除非我们先化解了我们与弟兄之间所结下的一切不和或冒犯。\BibleRef{玛 5:22} 因为不带平安就近天主的平安，这是何等的行为？保留债务却求豁免？他怎能平息他的\Emph{圣父}的怒气，而他自己却对自己的\Emph{弟兄}发怒，何况从一开始\Quote{一切愤怒}就被禁止？因为连若瑟在打发他的弟兄们去请他们的父亲时，也说过：\Quote{你们在路上不要生气。}他当时是在警告\Emph{我们}（因为别处我们称为\OriginalTerm{道路}{the Way}的纪律），当我们处于祈祷的\Quote{路上}时，不可带着怒气去见\Quote{圣父}。之后，主\Quote{扩大了法律}，公开将禁止对弟兄发怒添加到禁止杀人之上。\BibleRef{玛 5:21} 祂甚至不许以恶言发泄怒气。即使我们\Emph{必须}发怒，我们的怒气也不可留到日落之后，如宗徒所劝诫的。\BibleRef{厄 4:26} 但是，要么整天不祈祷，因为拒绝向弟兄赔补；要么因坚持发怒而失去祈祷，这何等鲁莽？\par

% structure: p0024 source=h2 target=section reason=relative-heading-rank; base=h2

\section{第十二章 同样必须免于一切心灵不安}

祈祷的操练不仅要从愤怒中，更要从\Emph{一切}心灵的不安中得到释放，要发自于与它所呈献的圣神相应的精神。因为污秽的精神不能被圣神所接纳，\BibleRef{厄 4:30} 忧伤的不能被喜乐的精神接纳，受束缚的不能被自由的精神接纳。没有人接待他的对手：除了他的同伴，没有人能进入。\par

% structure: p0026 source=h2 target=section reason=relative-heading-rank; base=h2

\section{第十三章 论洗手}

但是，带着确实洗净的手却怀着污秽的心灵去祈祷，这有何理由呢？——因为我们的手本身也需要属灵的洁净，使它们能\Quote{举起圣洁的手}\BibleRef{弟前 2:8}，远离虚假、谋杀、残暴、毒害、偶像崇拜以及所有其他由心灵孕育、通过手的操作产生的污点。这些才是真正的洁净；而不是那些大多数人所迷信地谨慎的事，每次祈祷都用水，即使刚从全身沐浴出来也是如此。当我仔细探究这一习俗并寻求其原因时，我发现它是一种纪念行为，关系到我们主的被出卖。\Emph{然而，我们向主祈祷}，我们并不\Emph{出卖}祂；相反，我们甚至应当反对那出卖者的榜样，并因此不洗手。除非在与人交往中沾染了污秽是洗手的一种良心原因，否则手本来是足够洁净的，因为我们已经与我们的整个身体一起在基督内一次洗净了。\par

% structure: p0028 source=h2 target=section reason=relative-heading-rank; base=h2

\section{第十四章 呼语}

尽管以色列每天洗净全身，却从未洁净。他的\Emph{手}无论如何总是污秽的，永远染着先知和主自己的血；因此，作为与他们父辈的罪行有份的世代罪犯，他们甚至不敢向主举手，唯恐某个依撒意亚呼喊\BibleRef{依 1:15}，唯恐基督全然战栗。然而，我们不仅举手，甚至伸开双手；并且，以主的受难为榜样，即使在祈祷中我们也向基督宣认。\par

% structure: p0030 source=h2 target=section reason=relative-heading-rank; base=h2

\section{第十五章 论脱去外衣}

但是既然我们已经触及了一个空泛礼规的特殊点，也不妨同样在其他点上打上我们的标记，因为这些点完全可以被指责为虚浮；也就是说，如果它们是在没有主或宗徒的任何诫命权威下被遵守的。因为这类事不属于宗教，而属于迷信，是人为的、强制的、出于好奇而非理性的仪式；无论如何都应当加以约束，即使是基于这一点：它们使我们与外邦人同列。例如，有些人习惯脱去外衣祈祷，因为外邦人就是这样接近他们的偶像；当然，如果这种实践是适宜的，那么教导有关祈祷服饰的宗徒们就会将其包含\Emph{在他们的训导中}，除非有人认为保禄是在祈祷时把他的外衣留在了卡尔颇那里！\BibleRef{弟后 4:13} 天主当然不会听穿外衣的恳求者，而祂明明听到了巴比伦王火窑中的三个圣人在穿着裤子和头巾祈祷。\par

% structure: p0032 source=h2 target=section reason=relative-heading-rank; base=h2

\section{第十六章 论祈祷后坐下}

再者，有些人祈祷结束后有坐下的习惯，我看不出有什么理由，除非是出于孩童般的借口。因为，如果那位\Person{赫耳玛斯}——他的作品通常冠以\WorkTitle{牧者}之名——在结束祈祷后，不是坐在床上，而是做了别的事，难道我们也要将此奉为规矩吗？当然不。为什么？即使经上说：\Quote{我祈祷完毕，坐在床上}，这句话也只是为了叙事的顺序，而非作为纪律的模本。否则，我们只能在有床的地方祈祷了！不，无论谁坐在\Emph{椅子}上或\Emph{长凳}上，都将与那著作相悖。再者，因为外邦人在敬拜他们的小神像后也坐着，因此这种做法在我们身上也该受责备，因为它见于偶像崇拜中。此外，还有\Emph{不敬}的罪责——即使外邦人若有理智也能明白。如果，一方面，在你最敬畏和尊崇的那一位眼前、面对面坐着是不敬的；那么，另一方面，在永生天主的眼前，当祈祷的天使仍\Emph{站在旁边}时，这样做更是\Emph{最}不虔诚的——除非我们在责备天主，说祈祷使我们疲倦了！\par

% structure: p0034 source=h2 target=section reason=relative-heading-rank; base=h2

\section{第十七章　论举手的姿势}

但是，当我们以谦逊和谦卑祈祷时，我们更将我们的祈祷托付给天主，双手不要举得过高，而是适度而合宜地举起；面容也不要过于张扬地抬起。因为那位在恳求中和面容上都以谦卑和沮丧祈祷的税吏，比那无耻的\Person{法利塞人}\BibleQuote{更为成义}地回去了。\BibleRef{路 18:9}同样，我们声音的音量也该压低；否则，如果我们因吵闹而被听见，我们需要多大的嗓门！但天主不是\Emph{声音}的听众，而是\Emph{心}灵的听众，正如祂是心灵的监察者。皮提亚神谕的恶魔说：\par

\Quote{我明白那沉默的，清楚听见那无声的。}\par

天主的耳朵等待声音吗？那么，\Person{约纳}的祈祷如何能从鲸鱼腹的深处，穿过如此巨兽的内脏，从深渊中，穿过如此浩瀚的海水，上达天堂？那些祈祷过于大声的人，除了打扰邻居外，能得到什么优越的好处？不，他们使自己的祈求可闻，这与在公开场合祈祷相比，犯的错误难道更少吗？\par

% structure: p0038 source=h2 target=section reason=relative-heading-rank; base=h2

\section{第十八章　论平安礼}

另一种习俗现已盛行。那些守斋的人在祈祷后，会免去平安礼——这平安礼是祈祷的印记。但是，什么时候比在举行某种宗教虔敬之时，我们的祈祷更具悦纳地上升，更应与弟兄们缔结平安呢？这样他们可以参与我们的虔敬，从而被软化，以便与他们的弟兄处理他们自己的平安。什么祈祷若离开了\Quote{圣吻}能算完整？平安妨碍谁为他主服务？什么样的祭献是使人离开时没有平安的？无论我们的祈祷是什么，它都不会比遵守那命令我们隐藏斋戒的诫命更好；\BibleRef{玛 6:16}因为\Emph{现在}，由于免去平安礼，我们被人知道是在守斋。但即使有某种理由\Emph{实行这做法}，仍要避免触犯这条诫命，你也许可以把\Quote{平安}推迟到\Emph{家中}，在那里你的斋戒不可能完全保密。但无论你在何处能隐藏你的虔敬行为，你应当记住这诫命：这样你可以在外满足纪律的要求，在家满足习俗的要求。同样，在逾越节那天，当守斋的宗教虔敬是普遍的，且可以说是公开的，我们正当地免去平安礼，不介意隐藏任何我们与众人共同做的事。\par

% structure: p0040 source=h2 target=section reason=relative-heading-rank; base=h2

\section{第十九章　论站斋}

同样，关于站斋的日子，多数人认为他们不应参加祭献的祈祷，理由是站斋必须通过领受主的身体而终止。那么，圣体圣事是取消还是更使一种献身于天主的服务与天主结合？你的\Emph{站斋}若同时\Emph{站立}在主的\Emph{祭台}前，岂不更庄严？当主的身体被领受并保留时，每个要点都得到保障：既参与了祭献，也履行了义务。如果\Quote{站斋}一名取自军营生活的榜样——因为我们也是天主的军队——当然，军营中没有喜乐或悲伤的歌唱会废除士兵的\Quote{站岗}：因为喜乐会更情愿地执行纪律，悲伤则更谨慎。\par

% structure: p0042 source=h2 target=section reason=relative-heading-rank; base=h2

\section{第二十章　论妇女的服饰}

然而，关于妇女的服饰，由于习俗的多样性，迫使卑微的我们，确实在至圣的宗徒之后，冒昧地加以讨论——除非我们按照宗徒来讨论，那就不是冒昧了。关于服饰和妆饰的谦逊，\Person{彼得}的指示\BibleRef{伯前 3:1}同样明确，他同\Person{保罗}以同一张口，因为怀着同一圣神，斥责衣着的荣耀、黄金的骄傲和头发的妖艳编梳。\par

% structure: p0044 source=h2 target=section reason=relative-heading-rank; base=h2

\section{第二十一章　论童贞女}

但那个在各个教会中众说纷纭的问题，即童贞女是否应当蒙头，必须加以讨论。因为，那些允许童贞女免于蒙头的人，似乎基于这一点：宗徒没有明确指明\Quote{童贞女}，而是指明\Quote{妇女}\BibleRef{格前 11:5}应\Quote{蒙头}；也没有泛指性别说\Quote{女性}，而是说了一个性别的\Emph{种类}，即\Quote{妇女}：因为如果他提到性别说\Quote{女性}，他就会为\Emph{每个}妇女划定绝对的界限；但他提到性别的一个种类时，就通过沉默将另一个种类区分开来。因为，他们说，他本可以特别提到\Quote{童贞女}，或者一般地用概括性词语\Quote{女性}。\par

% structure: p0046 source=h2 target=section reason=relative-heading-rank; base=h2

\section{第二十二章 对前述论点的答复}

凡作此让步者，当思这词本身的性质——自圣经最初记载起，\Quote{女人}一词有何含义？他们发现那是\Emph{性别的名称}，而非\Emph{性别的类别}：即天主在厄娃尚未认识男人之时，便赐予她\Quote{女人}和\Quote{女性}的称号——（\Quote{女性}泛指整个性别；\Quote{女人}则标志性别的某一类）。因此，既然当时尚未出嫁的厄娃被称为\Quote{女人}，这词便也成了童贞女的通称。宗徒——当然，是受同一圣神的指引，正如所有圣经，包括《创世记》在内，都是受那同一圣神默示而写成——在书写\Quote{女人}时用了同一个词，这词因未出嫁的厄娃之例，也能适用于\Quote{童贞女}，这并不奇怪。事实上，其余一切经文与此一致。因为正因他没有特别\Emph{指名}\Quote{童贞女}（如他在另一处教导婚姻问题时那样），便充分表明他的话是针对\Quote{所有}女人及\Quote{整个}性别而言；既然他压根没有指名童贞女，那么在\Quote{童贞女}与\Quote{其他女人}之间便没有区分。因为那在别处——即在区别需要之处——记得作区分的人（而且他以各自适当的名称来指称每一类别），在他不作区分（即不指名每一类别）之处，便希望人们明白并无差别。况且，在宗徒写信所用的希腊语中，通常说\Quote{女人}而非\Quote{女性}；即\OriginalTerm{女人}{\GreekText{γυναῖκας}}（\Emph{gunaikas}）而非\OriginalTerm{女性}{\GreekText{θηλείας}}（\Emph{theleias}）？因此，若那词（按翻译相当于\Quote{女性}（\Emph{femina}））常被用作性别的名称，那么他说\OriginalTerm{女人}{\GreekText{γυναῖκα}}时，便是以性别名称指称\Emph{性别}；而在\Emph{性别}中，连\Emph{童贞女}也包括在内。况且，声明是明白的：\Quote{\Emph{凡}女人祈祷或说预言，若不蒙头，就是羞辱自己的头。}\BibleRef{格前 11:5}\Quote{凡女人}若非\Quote{各种年龄、各种身份、各种状况的女人}，又是什么呢？他说\Quote{凡}，就丝毫不排除女人，正如他也不排除男人\Quote{不}蒙头一样；因为正如他所说：\Quote{\Emph{凡}男人}\BibleRef{格前 11:4}。因此，如同在男性中，以\Quote{男人}之名，连\Quote{少年}也被禁止蒙头；同样，在女性中，以\Quote{女人}之名，连\Quote{童贞女}也被命令蒙头。在每一性别中，让年幼者遵循年长者的规矩；否则，若女童贞女不被蒙头（因为她们未被指名），也请让男\Quote{童贞男}蒙头罢。让\Quote{男人}与\Quote{少年}有所区别，如果\Quote{女人}与\Quote{童贞女}也有所区别的话。因为他说女人必须蒙头，确实是\Quote{为了天使}\BibleRef{格前 11:10}，因为天使曾因\Quote{人的女儿们}而悖逆天主。那么，谁能主张唯有\Quote{\Emph{女人}}\Emph{独自}——即那些已婚且失去童贞者——才是天使贪恋的对象？除非\Quote{童贞女}不可能容貌出众、招引爱慕者？不，我们且看她们所贪恋的是否唯独\Emph{童贞女}？因为圣经说\Quote{\Emph{人的女儿们}}\BibleRef{创 4:2}，它本可平等地说\Quote{人的妻子们}或\Quote{女性}。同样，当它说：\Quote{他们就娶来为妻}\BibleRef{创 6:2}，这基于这样的理由：当然，那些被\Quote{娶为妻}的都是尚未有此名份的人。但对于那些并非如此的人，它会以不同方式表达。因此，（那些被指名的人）既无\Quote{寡居}之名，也无\Quote{童贞}之名。保禄通过泛指性别，完全将\Quote{女儿们}与各类别混入共性之中。再者，当他说\Quote{本性本身}\BibleRef{格前 11:14}——这本性已赐给女人头发作为遮盖和装饰——\Quote{教导蒙头是女人的本分}，难道童贞女就没有被赐予同样的遮盖和同样的头饰？若女人剪发是\Quote{可羞耻的}，对童贞女也同样如此。因此，对她们——她们被赋予了同一\Quote{头}的\Emph{律法}——也要求同一\Quote{头}的\Emph{规矩}（这规矩甚至延伸至那些因童年而受保护的童贞女，因为从一开始\Emph{童贞女}就被称为\Quote{女性}）。简言之，这习俗连\Emph{以色列}也遵守；但若\Emph{以色列}不遵守，则\Emph{我们}的律法（已被扩展和补充）将为自己辩护这附加的规定；且容许它也将面纱强加于童贞女。在\Emph{我们}这时代，让那不知自己性别的年龄保留单纯的自由。因为厄娃和亚当，当他们变得\Quote{有智慧}时\BibleRef{创 3:6}，便立即遮覆了他们所认识到的。无论如何，对于那些女子气已改变（为成熟）者，她们的年龄当记得对自然和规矩的本分；因为她们在人格和职责上都已被转入\Quote{女人}的行列。从她能结婚之时起，便不再是\Quote{童贞女}；因为在此时，年龄已嫁给了自己的丈夫，即时间。\Quote{但某位\Emph{童贞女}已献身于天主。从那一刻起，她既改变发式，也将全身装束改作\InnerQuote{女人}的样式。}那么，让她完全保持那身份，并履行\Quote{童贞女}的全部职责：她为天主所隐藏的，让她完全遮盖。将天主的恩宠在我们内所成就的，唯独托付于天主的认识，免得我们从人那里得到那我们从天主所望的赏报。为什么你在天主面前裸露你在人前所遮盖的？你会在公开场合比在教会中更矜持吗？若\Emph{你的献身}是天主的恩宠，而你已领受了它，\Quote{你为什么夸耀，}他说，\Quote{好像没有领受呢？}\BibleRef{格前 4:7}为何你因炫耀自己而论断别人？难道你因夸耀而引他人向善？不，连你自己也有丧失的危险，如果你夸耀的话；且你也将他人引向同样的危险！凡出于爱夸耀而采取的，容易败坏。童贞女，若你是童贞女，就当蒙头，因为你应当羞愧。若你是童贞女，就躲避众人的眼目。勿让人惊奇你的面容，勿让人察觉你的虚假。你若蒙头，倒是伪装已婚身份的好；不，你并非\Quote{伪装}——因为你已与基督成婚：你已向祂献上你的身体；行事要合你配偶的规矩。若祂命令他人的新娘蒙头，祂自己的新娘当然更当如此。\Quote{但各人不可认为前人的规矩可被推翻。}许多人屈服于他人的习俗，放弃自己的判断和一贯立场。即便童贞女\Quote{不被强迫}蒙头，那些\Quote{自愿}蒙头的也不应被禁止；她们同样不能否认自己是童贞女，却满足于在天主前良心的安稳，宁愿损害自己的名声。至于那些已订婚的，我可以以\Quote{超乎我微小度量}的坚定宣称并证明，她们从那天起——即当她们因亲吻和握手而初次因男子的身体接触而战栗时——就应当蒙头。因为在她们身上，一切都已预许：年龄因成熟，肉身因年龄，心灵因意识，谦逊因亲吻的经验，希望因期待，心智因意志。勒贝加对我们而言足为典范，她当未婚夫被指出时，仅因认出他便蒙上头纱为婚姻\BibleRef{创 24:64}。\par

% structure: p0048 source=h2 target=section reason=relative-heading-rank; base=h2

\section{第二十三章  论跪拜}

在\Emph{跪拜}事上，祈祷亦因若干人在安息日不跪拜而有不同的实践；且因这分歧尤其在教会面前受审，主将赐予恩宠，使异议者或让步，或在不冒犯他人的情况下坚持己见。然而我们（正如我们所领受的）只应在主复活之日，不仅避免跪拜，还应避免一切忧虑的姿态和职责；甚至推迟我们的俗务，免得给魔鬼留地步。\BibleRef{弗 4:27}同样，在五旬节期间也应如此；我们以相同的欢跃礼仪来区分这时期。但谁又会在每日至少第一次祈祷、在晨光中进入祈祷时，不向天主俯伏呢？此外，在斋戒和站岗时，若无跪拜及其余惯常的谦卑标记，便不应祈祷；因为那时我们不单是\Emph{祈祷}，更是\Emph{祈求}，并向主天主做补赎。至于祈祷的\Emph{时间}，则毫无规定，除非明确地说\Quote{要时时处处祈祷}。\par

% structure: p0050 source=h2 target=section reason=relative-heading-rank; base=h2

\section{第二十四章  论祈祷的地方}

但如何\Quote{在每一个地方}祈祷，既然我们被禁止在公共场所祈祷？他指的是每一个机会或甚至必要使之适宜的地方：因为宗徒们（在狱中，当着囚犯的面，\Quote{开始祈祷并歌颂天主}）所做的不被认为违反诫命；保禄在船上，当着众人\Quote{感谢天主}也不违反。\par

% structure: p0052 source=h2 target=section reason=relative-heading-rank; base=h2

\section{第二十五章  论祈祷的时间}

至于\Emph{时间}，遵守某些时辰的外在习惯并非无益——我指的是那些标示一天时段的通常时辰：第三时辰、第六时辰、第九时辰，我们在圣经中可发现这些时辰比其他时辰更为隆重。圣神初次倾注于聚集的门徒身上是在\Quote{第三时辰}。伯多禄在他看到普世团体异象的那一天，在那小船中，为了祈祷而登上\Emph{房屋}的更高处，是在\Quote{第六时辰}。\BibleRef{宗 10:9}同一宗徒与若望进入圣殿是在\Quote{第九时辰}，那时他使瘫痪者恢复健康。虽然这些\Emph{实行}并无任何\Emph{诫命}要求遵守，但确立某种明确的习惯仍是一件好事，它既能加强祈祷的劝勉，又能像法律一样将我们从俗务中拉出来履行这职责；这样我们至少一日祈祷三次，有如达尼尔所遵守的（当然符合以色列的纪律），因为我们欠三位——圣父、圣子、圣神的债；当然，此外还有因晨光和夜晚来临而应做的常规祈祷，这些无需提醒。而且，信徒在进食和沐浴之前，也必须先祈祷；因为精神的滋养和食粮应先于肉身的，天上的应先于地上的。\par

% structure: p0054 source=h2 target=section reason=relative-heading-rank; base=h2

\section{第二十六章  论兄弟的告别}

你不可不祈祷就打发进入你家的兄弟离开。——\Quote{你看见了，}\Emph{圣经}说，\Quote{一个兄弟？就是看见了你的主；}——尤其是\Quote{一个陌生人}，免得他可能是\Quote{一位天使}。但反之，当你自己受到弟兄接待时，不可使地上的饮食先于天上的，因为你的信德将立即受到审判。否则，你如何——按照诫命\BibleRef{路 10:5}——说\Quote{愿这\Emph{家}平安}，除非你与家中的那些人交换平安？\par

% structure: p0056 source=h2 target=section reason=relative-heading-rank; base=h2

\section{第二十七章  论附加圣咏}

更勤于祈祷的人，通常在祈祷中附加\Quote{阿肋路亚}及此类圣咏，会众在结尾回应。当然，每项旨在赞扬和光荣天主、齐心将更丰富的祈祷作为精选祭品献上的制度，都是卓越的。\par

% structure: p0058 source=h2 target=section reason=relative-heading-rank; base=h2

\section{第二十八章  论作为属神祭品的祈祷}

因为这就是那属神的祭品\BibleRef{伯前 2:5}，它取代了昔日的牺牲。他说：\Quote{你们的众多祭品对我有何用？我饱享了公羊的全燔祭，不喜爱公羊的脂肪，也不喜爱公牛和山羊的血。谁向你们索要这些呢？}那么，天主\Emph{确曾}索要什么，福音教导了。他说：\Quote{时候将到，那时真正朝拜的人将在心神和真理中朝拜父。因为天主是神，因此要求朝拜他的人也是如此。}\BibleRef{若 4:23}我们就是真正的朝拜者和真正的司祭，我们以心神祈祷，在心神中献上祈祷——这是天主所悦纳的合适祭品，是他所索要、所期盼的！这\Emph{祭品}，以全心奉献，以信德滋养，以真理培育，纯洁无瑕，贞洁无染，以爱德为冠冕，我们应以善行的行列、在圣咏与赞歌中，护送它到天主的祭台前，好从天主那里为我们求得一切。\par

% structure: p0060 source=h2 target=section reason=relative-heading-rank; base=h2

\section{第二十九章  论祈祷的能力}

因为那要求我们以\Quote{心神和真理？}祈祷的天主，又\Emph{何曾}拒绝过祈祷呢？我们所诵读、所听闻、所相信的，尽是祈祷大能的有力见证！\Emph{古时}的祈祷确曾使人脱离烈火\EditorialNote{达内尔 3}、脱离野兽\EditorialNote{达内尔 6}和饥荒；然而它那时尚未从基督领受自己的形式。但\Emph{基督徒}的祈祷其功效是何等丰盛！它不再将甘露天使安置在烈火中，不再给狮子戴上口套，不再将乡人的饼送给饥者\BibleRef{列下 4:42}；它没有分派的恩宠以消除一切痛苦之感；却以坚忍供养受苦、感受和悲伤的人：它借德行增扩恩宠，好让信德知道她从主那里获得了什么，并明白她为了天主的名而忍受了什么。然而在过去的日子里，祈祷也曾召来瘟疫、驱散敌军、阻挡甘霖的滋养。然而如今，义人的祈祷平息天主的一切愤怒，为私仇宿敌守望，为迫害者祈求。既然\Emph{它}知道如何强取天上的\Emph{甘霖}——（祈祷）\Emph{曾经}能求得\Emph{烈火}，这又有什么可惊奇的？唯独祈祷能战胜天主。但基督愿意祈祷不行恶事：祂已将祈祷的一切能力都赐予善工。因此，祈祷只懂得：将逝者的灵魂从死亡之道上召回，转变软弱者，治愈病者，驱魔，打开监门，解开无辜者的锁链。同样，它洗除过犯，击退诱惑，扑灭迫害，安慰心怯者，鼓舞心高者，护送旅人，平息波涛，令盗贼惊骇，养育穷人，管理富人，扶起跌倒者，扶持将倾者，坚固站立者。祈祷是信德的城墙：是对那四面窥伺我们的仇敌的武器和投枪。因此，我们绝不可赤手空拳而行。白日，我们当谨记\OriginalTerm{站（斋戒日）}{Station}；夜间，谨记守夜。在祈祷的武器之下，我们守护我们\OriginalTerm{统帅}{General}的旗帜；我们在祈祷中等待天使的号角。同样，众天使也都祈祷；一切受造物都祈祷；牲畜和野兽祈祷并屈膝；当它们走出巢穴和洞穴时，它们向上仰望天庭，嘴巴并非徒然张开，以自己特有的方式使气息颤动。不仅如此，鸟儿也从巢中飞起，向天高举自己，并以翅膀展开十字架形代替双手，发出某种似乎祈祷的声音。那么，关于祈祷的职分还能说什么呢？连主自己也曾祈祷；愿光荣与德能归于祂，至于无穷之世！\par

\SourceRecord{Translated by S. Thelwall.；Ante-Nicene Fathers；Vol. 3.；Edited by Alexander Roberts, James Donaldson, and A. Cleveland Coxe.；Buffalo, NY: Christian Literature Publishing Co.,；1885.；<http://www.newadvent.org/fathers/0322.htm>.}
